% Source-derived chapter 01: Introduction
% Original source: pixtons-formula-abel-jacobi-picard-stack
\chapter{Introduction}
\label{pixtons-formula-abel-jacobi-picard-stack-chapter-01}
\source{pixtons-formula-abel-jacobi-picard-stack}

\begin{remark}[Source section]
\label{pixtons-formula-abel-jacobi-picard-stack-chapter-01-source}
\source{pixtons-formula-abel-jacobi-picard-stack}
\source{pixtons-formula-abel-jacobi-picard-stack}
This chapter reproduces the corresponding section of the retrieved
LaTeX source, including its definitions, statements, examples, and
proofs. It has not yet been translated into Lean.
\notready
\end{remark}

% The source section title was: Introduction
\parskip=5pt


\subsection{Double ramification cycles}
\label{Ssec:DRclassic}
\source{pixtons-formula-abel-jacobi-picard-stack}

Let $A = (a_1, \dots, a_n)$ be a vector of $n$ integers satisfying
%{\footnote{In the twisted case,
%the vanishing sum condition will be removed.}}
$$\sum_{i=1}^n a_i = 0\,. $$ In the moduli space $\cM_{g,n}$ of 
nonsingular curves of genus $g$ with $n$ marked points,
 consider the substack defined by the following classical condition: 
\begin{equation}\label{g445}
\source{pixtons-formula-abel-jacobi-picard-stack}
\left\{ (C, p_1, \dots, p_n) \in \cM_{g,n} \; \rule[-0.8em]{0.05em}{2em} \; 
\cO_C\Big(\sum_{i=1}^n a_i p_i\Big) \simeq \cO_C \right\}\, . 
\end{equation}
%nd it is natural to ask what is the homology class of its closure or of its natural compactification in $\oM_{g,n}$.
From the point of view of relative Gromov-Witten theory, 
the most natural compactification of the substack \eqref{g445} 
is the space $\oM^{\sim}_{g,A}$ of stable maps to
{\em rubber}\,: stable maps to
 $\CP^1$ relative to $0$ and $\infty$ modulo the $\C^*$-action on $\CP^1$.
The rubber moduli space carries a natural virtual fundamental class $\left[\oM^{\sim}_{g,A}\right]^\virt$ of (complex)
dimension $2g-3+n$. The pushforward 
via the canonical morphism
$$ \epsilon:\oM^{\sim}_{g,A} \rightarrow  \oM_{g,n}$$
is the {\em double ramification cycle} 
$$\epsilon_*\left[\oM^{\sim}_{g,A}\right]^\virt\, =\, \mathsf{DR}_{g,A}\ \in 
{\mathsf{CH}}_{2g-3+n}(\overline{\cM}_{g,n})\,. $$


Double ramification cycles have been studied intensively for the 
past two decades.
Examples of early results can be found in \cite{BSSZ,Cavalieri2011Polynomial-fami,cj,Faber2005Relative-maps-a,Grushevsky2012The-double-rami,Hain2013Normal-function, Marcus2013Stable-maps-to-}.
A complete formula was conjectured by
Pixton in 2014 and proven in \cite{Janda2016Double-ramifica}. 
For subsequent 
study and applications, see \cite{Bae,BGR,CGJZ,FanWu,Farkas2016The-moduli-spac,Holmes2017Extending-the-d,
Holmes2017Jacobian-extens,
Holmes2017Multiplicativit,MPS20,ObPix,Pix18,Schmitt2016Dimension-theor,
Tseng2,Tseng1}. Essential for our work is the formula for double ramification
cycles for target varieties in \cite{Janda2018Double-ramifica}.

We refer the
reader to \cite[Section 0]{Janda2016Double-ramifica} and \cite[Section 5]{RPSLC} for introductions to the subject. For a classical perspective from the point of view
of Abel-Jacobi theory, see 
\cite{Holmes2017Extending-the-d}.



\subsection{Twisted double ramification cycles}
We develop here a theory which extends the study of double ramification
cycles from the moduli space of stable  curves $\overline{\cM}_{g,n}$ to the Picard stack of
curves with line bundles $\Picabs_{g,n}$.
An object of $\Picabs_{g,n}$ over $\mathcal{S}$ is
a flat family $$\pi:\mathcal{C} \rightarrow \mathcal{S}$$ of prestable{\footnote{A prestable $n$-pointed
curve is a connected nodal curve with markings at distinct nonsingular points.
For the entire paper, we avoid the $(g,n)=(1,0)$ case because of 
non-affine stabilizers.}}
    $n$-pointed genus $g$ curves together with a line bundle
    $$\mathcal{L} \rightarrow \mathcal{C}\,. $$
The Picard stack $\Picabs_{g,n}$ is an algebraic (Artin) stack
which is locally of finite type, see \cref{sec:pic_stacks} for a treatment
of foundational issues. 


Since the degree of a line bundle is constant in flat families, there is a
disjoint union
$$\Picabs_{g,n} =\bigcup_{d\in \mathbb{Z}}  \Picabs_{g,n,d}\, ,$$
where $\Picabs_{g,n,d}$ is the Picard stack of curves with degree $d$
line bundles.
Let $$A = (a_1, \dots, a_n)\,,\ \ \ \ \  \sum_{i=1}^n a_i=d \,$$
be a vector of integers.
The first result of the paper is the construction of a 
{\em universal twisted double ramification
cycle} in the operational Chow theory{\footnote{All
Chow theories in the paper will be taken with $\mathbb{Q}$-coefficients.
}} of $\Picabs_{g,n,d}$,
$$\mathsf{DR}^{\mathsf{op}}_{g,A} \in {\mathsf{CH}}^g_{\mathsf{op}}(\Picabs_{g,n,d})\, .$$

The geometric intuition behind the construction is simple.
Let 
$$\pi:\mathcal{C}\rightarrow \mathcal{S}\, ,
\ \ \ \ p_1,\ldots,p_n: \mathcal{S} \rightarrow \mathcal{C}\, ,
\ \ \ \ \mathcal{L}\rightarrow \mathcal{C}$$
be an object of $\Picabs_{g,n,d}$.
The class
$\mathsf{DR}^{\mathsf{op}}_{g,A}$ should operate as the 
locus in the base $\mathcal{S}$ heuristically determined by the condition
$$
\cO_C\Big(\sum_{i=1}^n a_i p_i\Big) \simeq \mathcal{L}|_{C} \,.$$

%To make the above idea precise, we do {\em not} use the virtual class of the moduli space of stable
%maps in Gromov-Witten theory, but rather an alternative approach via log geometry based closely on the stack of %tropical divisors constructed in \cite{Marcus2017Logarithmic-com}. Our construction is presented in Section %\ref{sec:log_def}. We also give a naive interpretation in terms of the closure of the unit section in  $\Picabs$ in the spirit of \cite{Holmes2017Jacobian-extens} (Section \ref{sec:naive_def}), and an interpretation in terms of resolving the indeterminacies of the classical Abel-Jacobi map in the manner of \cite{Holmes2017Extending-the-d} (Section \ref{Sect:AJextension}, where we also describe a b-Chow version in the spirit of \cite{Holmes2017Multiplicativit}). 



To make the above idea precise, we do {\em not} use the virtual class of the moduli space of stable
maps in Gromov-Witten theory, but rather an alternative approach by
partially resolving the classical Abel-Jacobi map.
The method follows the path 
of \cite{Holmes2017Extending-the-d, Holmes2017Multiplicativit}
and
may be viewed  as a universal Abel-Jacobi construction
over the Picard stack. 
Log geometry based on the stack of tropical divisors constructed in \cite{Marcus2017Logarithmic-com}
plays a crucial role.
Our construction is presented in Section \ref{Sect:AJextension}.
%\jocomment{This is close to what we planned to discuss in Section \ref{Sect:AJextension}?}

The basic compatibility of our new operational class $$\mathsf{DR}^{\mathsf{op}}_{g,A} \in {\mathsf{CH}}^g_{\mathsf{op}}(\Picabs_{g,n,d})$$ with the
standard double ramification cycle is as follows.
Let 
$$A = (a_1, \dots, a_n)\,,\ \ \ \ \  \sum_{i=1}^n a_i=0 \,,$$
be given. 
The universal data 
\begin{equation}\label{fftt66677}
\source{pixtons-formula-abel-jacobi-picard-stack}
\pi: \mathcal{C}_{g,n} \rightarrow \overline{\cM}_{g,n}\, , \ \ \ \ \cO \rightarrow \mathcal{C}_{g,n}
\end{equation}
determine a map $\phi_{\ca O} : \overline{\cM}_{g,n} \to \Picabs_{g,n,0}$.
%Consider the universal family of over the moduli space of
%stable curves
%\begin{equation}\label{fftt55}
%\pi: \mathcal{C}_{g,n} \rightarrow \overline{\cM}_{g,n}
%\end{equation}
%equipped with the trivial line bundle
%\begin{equation}\label{fftt66}
%%\cO_{\mathcal{C}_{g,n}} \rightarrow {\mathcal{C}}_{g,n}\, .
%\end{equation}
%The data \eqref{fftt66677} and \eqref{fftt66} together determine an object of
%$\Picabs_{g,n,0}$.
The action of $\mathsf{DR}^{\mathsf{op}}_{g,A}$ on the fundamental class
of $\overline{\mathcal{M}}_{g,n}$ corresponding to the family \eqref{fftt66677} then equals the previously defined
double ramification cycle
$$\mathsf{DR}^{\mathsf{op}}_{g,A}(\phi_{\ca O})\Big( [\overline{\mathcal{M}}_{g,n}]\Big)= \mathsf{DR}_{g,A}
\in {\mathsf{CH}}_{2g-3+n}(\overline{\mathcal{M}}_{g,n})
\, .$$




More generally, 
for a vector $A = (a_1, \dots, a_n)$ of integers satisfying
$$\sum_{i=1}^n a_i = k(2g-2)\,, $$ 
canonically twisted double ramification cycles,
$$\mathsf{DR}_{g,A, \omega^k}
\in {\mathsf{CH}}_{2g-3+n}(\overline{\mathcal{M}}_{g,n})\, ,$$
related to the classical loci
\begin{equation*}
\left\{ (C, p_1, \dots, p_n) \in \cM_{g,n} \; \rule[-0.8em]{0.05em}{2em} \; 
\cO_C\Big(\sum_{i=1}^n a_i p_i\Big) \simeq  \omega^k_C \right\},\,  
\end{equation*}
have been constructed in
\cite{Guere2016A-generalizatio} for $k=1$ and in \cite{Holmes2017Extending-the-d,Holmes2017Jacobian-extens,Marcus2017Logarithmic-com} for all $k\geq 1$.
The universal data 
\begin{equation}\label{fftt66672}
\source{pixtons-formula-abel-jacobi-picard-stack}
\pi: \mathcal{C}_{g,n} \rightarrow \overline{\cM}_{g,n}\, , \ \ \ \ \omega_\pi^k \rightarrow \mathcal{C}_{g,n}
\end{equation}
determine a map 
$\phi_{\omega_\pi^k} : \overline{\cM}_{g,n} \to \Picabs_{g,n,k(2g-2)}$. 
Here,
$\omega_\pi$ is the relative dualizing sheaf
of the morphism $\pi$.

The action of $\mathsf{DR}^{\mathsf{op}}_{g,A}$ on the fundamental class
of $\overline{\mathcal{M}}_{g,n}$ corresponding to the family \eqref{fftt66672} is compatible with the constructions of \cite{Guere2016A-generalizatio,Holmes2017Extending-the-d,Holmes2017Jacobian-extens,Marcus2017Logarithmic-com},
$$\mathsf{DR}^{\mathsf{op}}_{g,A} (\phi_{\omega_\pi^k})\Big( [\overline{\mathcal{M}}_{g,n}]\Big)= \mathsf{DR}_{g,A, \omega^k}
\in {\mathsf{CH}}_{2g-3+n}(\overline{\mathcal{M}}_{g,n})
\, $$
for all $k\geq 1$.

The above compatibilities of $\mathsf{DR}^{\mathsf{op}}_{g,A}$ with
the standard and canonically twisted double ramification cycles are proven in Section \ref{proof1111}.

\begin{theorem}
\source{pixtons-formula-abel-jacobi-picard-stack}
\notready \label{cccc}
\source{pixtons-formula-abel-jacobi-picard-stack}
Let $g\geq 0$ and $d\in \mathbb{Z}$.
Let $A=(a_1,\ldots,a_n)$ be a vector of integers satisfying
$$\sum_{i=1}^n a_i=d\, .$$
Logarithmic compactification of the Abel-Jacobi map yields a universal twisted double
ramification 
cycle
$$\mathsf{DR}^{\mathsf{op}}_{g,A} \in {\mathsf{CH}}^g_{\mathsf{op}}(\Picabs_{g,n,d})$$
which is compatible with the standard double ramification
cycle 
$$\mathsf{DR}_{g,A, \omega^k}\in \mathsf{CH}_{2g-3+n}(\overline{\cM}_{g,n})\, $$
in case $d=k(2g-2)$ for $k\geq 0$.
\end{theorem}

%\{Do we want to mention the targets version here also? I see there is a remark of Johannes here along these %lines commented out, so maybe this has already been discussed? Rahul: no, too much for the intro}

%\jocomment{This seems weaker than the following theorem.}
%\begin{theorem*}
%Let $g\geq 0$, then partial resolution of the Abel-Jacobi map yields a universal twisted double
%ramification 
%cycle
%$$\mathsf{DR}^{\mathsf{op}}_{g} \in {\mathsf{CH}}^g_{\mathsf{op}}(\Picabs_{g,0,0}).$$
%For $A=(a_1, \ldots, a_n) \in \mathbb{Z}^n$ with $\sum_{i=1}^n a_i = k (2g-2)$ the pullback of this cycle by the map
%\[\overline{\cM}_{g,n} \to \Picabs_{g,0,0}, (C,p_1, \ldots, p_n) \mapsto \left(C, \omega_C^{\otimes k} \left( -\sum_{i=1}^n a_i p_i\right) \right),\]
%applied to the fundamental class $[\overline{\cM}_{g,n}]$ gives Holmes' definition of the Double ramification cycle \cite{Holmes2017Extending-the-d}. In particular, for $k=0$ this also recovers the Double ramification cycle defined using the space of rubber relative maps. \jocomment{Similarly, compatibility with DR cycle with targets.}
%\end{theorem*}


\subsection{Pixton's formula}

\subsubsection{Prestable graphs}\label{xvalsg}
\source{pixtons-formula-abel-jacobi-picard-stack}
We define the set $\G_{g,n}$ of {\em prestable graphs} as
follows. A prestable graph $\Gamma \in \G_{g,n}$ consists of the data
$$\Gamma=(\V\, ,\ \H\, ,\ \L\, , \ \mathrm{g}:\V \rarr \Z_{\geq 0}\, ,
\ v:\H\rarr \V\, , 
\ \iota : \H\rarr \H)$$
satisfying the properties:
\begin{enumerate}
\item[(i)] $\V$ is a vertex set with a genus function $\g:\V\to \Z_{\geq 0}$,
\item[(ii)] $\H$ is a half-edge set equipped with a 
vertex assignment $v:\H \to \V$ and an involution $\iota$,
\item[(iii)] $\E$, the edge set, is defined by the
2-cycles of $\iota$ in $\H$ (self-edges at vertices
are permitted),
\item[(iv)] $\L$, the set of legs, is defined by the fixed points of $\iota$ and is
placed in bijective correspondence with a set of $n$ markings,
\item[(v)] the pair $(\V,\E)$ defines a {\em connected} graph
satisfying the genus condition 
$$\nice \sum_{v \in \V} \g(v) + h^1(\Gamma) = g\, .$$
%\item[(vi)] for each vertex $v$, the stability condition holds: 
%$$2\g(v)-2+ \n(v) >0\, ,$$
%where $\n(v)$ is the valence of $\Gamma$ at $v$ including 
%both edges and legs.
\end{enumerate}
To emphasize $\Gamma$, the notation $\V(\Gamma)$, $\H(\Gamma)$, $\L(\Gamma)$, and $\E(\Gamma)$ will also be used for the vertex, half-edges, legs, and
edges of $\Gamma$. 

An isomorphism between $\Gamma$, $\Gamma' \in \G_{g,n}$ 
consists of bijections $\V \to \V'$ and $\H \to \H'$ respecting the
structures $\L$, $\mathrm{g}$, $v$, and $\iota$.
Let $\Aut(\Gamma)$ denote the automorphism group of $\Gamma$.

% An automorphism of $\Gamma \in \G_{g,n}$ 
% consists of automorphisms
% of the sets $\V$ and $\H$ which leave invariant the
% structures $\L$, $\mathrm{g}$, $v$, and $\iota$.
% Let $\Aut(\Gamma)$ denote the automorphism group of $\Gamma$.


While the set of isomorphism classes of prestable graphs is infinite, the set of isomorphism classes of prestable graphs with prescribed bounds on the number of edges is finite.
% subsets consisting of
% prestable graphs with prescribed bounds on the number of edges contain only finitely many isomorphism classes. 



Let $\mathfrak{M}_{g,n}$ be the algebraic (Artin) stack of prestable curves of genus $g$ with $n$ marked points.
A prestable graph $\Gamma$ determines an algebraic stack $\mathfrak{M}_\Gamma$ of 
curves
with degenerations forced by the graph,
$$
\mathfrak{M}_\Gamma = \prod_{v \in \V} \mathfrak{M}_{\g(v), \n(v)}\, $$
together with a 
canonical\footnote{To define the map, we choose an ordering on the half-edges at each vertex. } map
$$j_\Gamma : \mathfrak{M}_\Gamma \to \mathfrak{M}_{g,n}\, .$$
Since $\mathfrak{M}_{g,n}$ is smooth and
the morphism $j_\Gamma$ is proper, representable, and lci, we obtain an operational Chow class on the algebraic
stack of curves,
$$j_{\Gamma*}[\mathfrak{M}_\Gamma] \in \mathsf{CH}_{\mathsf{op}}^{|\E(\Gamma)|}(\mathfrak{M}_{g,n})\, .$$
Via the morphism of algebraic stacks,
$$\epsilon: \Picabs_{g,n,d} \rightarrow \mathfrak{M}_{g,n}\, ,$$
$j_{\Gamma*}[\mathfrak{M}_\Gamma]$ also defines an operational Chow class
on the Picard stack,
$$\epsilon^* j_{\Gamma*}[\mathfrak{M}_\Gamma] \in \mathsf{CH}_{\mathsf{op}}^{|\E(\Gamma)|}(\Picabs_{g,n,d})\, .$$

\subsubsection{Prestable graphs with degrees} \label{Sect:Prestgrwdeg}
\source{pixtons-formula-abel-jacobi-picard-stack}
We will require a refinement of the prestable graphs of Section \ref{xvalsg} which includes degrees
of line bundles. 

We define the set $\G_{g,n,d}$ of {\em prestable graphs of degree $d$} as
follows: $$\Gamma_\delta=(\Gamma,\delta) \in \G_{g,n,d}$$ consists of the data
\begin{enumerate}
    \item[$\bullet$] a prestable graph $\Gamma\in \mathsf{G}_{g,n}$,
    \item[$\bullet$] a function $\delta: \V \rightarrow \mathbb{Z}$
    satisfying the degree condition
$$\sum_{v\in V} \delta(v) = d\, .$$
\end{enumerate}
The function $\delta$ is  often called the {\em multidegree}.


An automorphism of $\Gamma_\delta \in \G_{g,n,d}$ 
consists of an automorphism of $\Gamma$
leaving $\delta$ invariant.
Let $\Aut(\Gamma_\delta)$ denote the automorphism group of $\Gamma_\delta$.


%For a prestable graph $\Gamma\in \mathsf{G}_{g,n}$ and a function $\delta:V \rightarrow \mathbb{Z}$
%which satisfies the degree condition (vi), let $\Gamma_\delta\in \mathsf{G}_{g,n,d}$ be the associated
%prestable graph of degree $d$.

For $\Gamma_\delta\in \mathsf{G}_{g,n,k}$, let $\mathfrak{M}_\Gamma$ be the algebraic stack of
curves defined in Section \ref{xvalsg} with respect to the underlying prestable graph $\Gamma$. 
Let
$\Picabs_{\Gamma_\delta}$ be the Picard stack,
$$\epsilon:\Picabs_{\Gamma_\delta} \rightarrow \mathfrak{M}_{\Gamma}\, ,$$
parameterizing curves with degenerations forced by $\Gamma$ and with  line bundles which
have 
degree $\delta(v)$ restriction to the components corresponding to the vertex $v\in \V$.
We have a
canonical map
$$j_{\Gamma_\delta} : \Picabs_{\Gamma_\delta} \to \Picabs_{g,n,d}\, .$$
Since $\Picabs_{g,n,d}$ is smooth and
the morphism $j_{\Gamma_\delta}$ is proper, representable, and lci,
we obtain an operational Chow class,
$$j_{\Gamma_\delta *}[\Picabs_{\Gamma_\delta}] \in \mathsf{CH}_{\mathsf{op}}^{|\E(\Gamma)|}(\Picabs_{g,n,d})\, .$$
As operational Chow classes, the following formula holds:
\begin{equation}\label{5599w}
\source{pixtons-formula-abel-jacobi-picard-stack}
\epsilon^*j_{\Gamma*}[\mathfrak{M}_\Gamma] = \sum_{\delta} j_{\Gamma_\delta *}[\Picabs_{\Gamma_\delta}]
\ \in \mathsf{CH}_{\mathsf{op}}^{|\E(\Gamma)|}(\Picabs_{g,n,d})
\, ,
\end{equation}
where the sum{\footnote{The sum is infinite, but only finitely many terms
are nonzero in any operation.}} is over all functions $\delta:\V \rightarrow \mathbb{Z}$ satisfying the degree condition.
Equivalently, we may write \eqref{5599w} as
\begin{equation*}
\frac{1}{|{\text{Aut}}(\Gamma)|}\, \epsilon^*j_{\Gamma*}[\mathfrak{M}_\Gamma] = 
\sum_{\Gamma_\delta\in \mathsf{G}_{g,n,d}} 
\frac{1}{|{\text{Aut}}(\Gamma_\delta)|}\, 
j_{\Gamma_\delta *}[\Picabs_{\Gamma_\delta}]\ 
\in \mathsf{CH}_{\mathsf{op}}^{|\E(\Gamma)|}(\Picabs_{g,n,d})
\, ,
\end{equation*}
where the sum on the right side is now over all isomorphism classes of
prestable graphs of degree $d$ with underlying prestable graph $\Gamma$.

\subsubsection{Tautological \texorpdfstring{$\psi$}{psi}, \texorpdfstring{$\xi$}{xi}, and \texorpdfstring{$\eta$}{eta} classes} \label{Ssec:PsiXiEta}
\source{pixtons-formula-abel-jacobi-picard-stack}

The universal curve $$\pi:\mathfrak{C}_{g,n} \to \Picabs_{g,n}$$ 
carries two natural line bundles: the relative dualizing sheaf $\omega_{\pi}$ 
and the universal line bundle
$$\mathfrak{L} \rightarrow \mathfrak{C}_{g,n}\, .$$
Let $p_i$ be the $i$th section of the universal curve, let 
$$\mathfrak{S}_i\subset \mathfrak{C}_{g,n}$$
be  the corresponding divisor, and let
$$\omega_\log = \omega_\pi\Big(\sum_{i=1}^n \mathfrak{S}_i\Big)$$ be the relative log-canonical line bundle 
with first Chern class $c_1 (\omega_\log)$. 
Let $$\xi = c_1(\mathfrak{L})$$ be the first Chern class of $\mathfrak{L}$.

\begin{definition}
\source{pixtons-formula-abel-jacobi-picard-stack}
\notready \label{Not:PsiXiEta}
\source{pixtons-formula-abel-jacobi-picard-stack}
The following operational classes on $\Picabs_{g,n}$ are obtained from the
universal structures:
\begin{itemize}
    \item $\psi_i = c_1(p_i^* \omega_\pi)\in \mathsf{CH}^1_{\mathsf{op}}(\Picabs_{g,n})\, $,
    \item $\xi_i = c_1(p_i^* \mathfrak{L})\in \mathsf{CH}^1_{\mathsf{op}}(\Picabs_{g,n})
    \, $,
    \item $\eta_{a,b} = \pi_*\left(c_1(\omega_\log)^a \xi^b\right)
    \in \mathsf{CH}^{a+b-1}_{\mathsf{op}}(\Picabs_{g,n})
    \, $.
\end{itemize}
\end{definition}
For simplicity in the formulas, 
we will use the
notation
$$\eta= \eta_{0,2} =\pi_*(\xi^2)\, .$$
The standard $\kappa$ classes are defined by the $\pi$ pushforwards
of powers of $c_1(\omega_\log)$, 
$$\eta_{a,0} = \kappa_{a-1}\, .$$

\begin{definition}
\source{pixtons-formula-abel-jacobi-picard-stack}
\notready
A {\em decorated prestable graph} $[\Gamma_\delta,\gamma]$ {\em of degree $d$} 
%of type $(g,n,\beta)$ 
is a prestable graph $\Gamma_\delta \in \G_{g,n,d}$ of degree $d$ together with the following
decoration data $\gamma$:
\begin{itemize}
\item each leg $i\in \L$ is decorated with a monomial $\psi_i^a \xi_i^b$,
\item each half-edge $h\in \H\setminus \L$ is 
decorated with a monomial $\psi^a_h$,
\item each edge $e\in \E$ is decorated with a monomial $\xi^a_e$,
\item each vertex in $\V$ is decorated with a monomial in the variables 
$\{\eta_{a,b}\}_{a+b \geq 2}$.
\end{itemize}
In all four cases, the monomial may be be trivial.
\end{definition}


Let $\D\G_{g,n,d}$ be the set of decorated prestable graphs of degree $d$.
To each decorated graph of degree $d$, 
$$[\Gamma_\delta,\gamma]\in \D\G_{g,n,d}\,,$$ we assign the 
operational class 
$$j_{\Gamma_\delta*}[\gamma] \in \mathsf{CH}^*_{\mathsf{op}}(\Picabs_{g,n,d})$$
obtained
via the proper representable morphism
$$j_{\Gamma_\delta} : \Picabs_{\Gamma_\delta} \rightarrow \Picabs_{g,n,d}$$
and the action of the decorations.

The action of decorations is described
as follows.
Given $\Gamma_\delta\in \mathsf{G}_{g,n,k}$, the stack $\Picabs_{\Gamma_\delta}$ admits a morphism{\footnote{The fibers of the map are torsors under the group $\mathbb{G}_m^{h^1(\Gamma)}$.}}
\[\Picabs_{\Gamma_\delta} \to \prod_{v \in \V(\Gamma_\delta)} \Picabs_{\g(v),\n(v),\delta(v)}\]
which sends a line bundle $\mathcal{L}$ on a prestable curve $\mathcal{C}$ to its restrictions on the various components,
$$ \mathcal{L}|_{\mathcal{C}_v}\, , \  \ \  v\in \V(\Gamma_\delta)\, . $$
For $v \in \V(\Gamma_\delta)$, we define the operational
class $\eta(v)$ on $\Picabs_{\Gamma_\delta}$ as the pullback of the operational class $\eta$ on the factor $\Picabs_{\g(v),\n(v),\delta(v)}$ above. The operational
classes $\psi$ at the markings and $\xi$ at the half-edges are
defined similarly.




\begin{definition}
\source{pixtons-formula-abel-jacobi-picard-stack}
\notready The {\em tautological  classes} in
 $\mathsf{CH}^*_{\mathsf{op}}(\Picabs_{g,n,d})$ consist of the 
 $\mathbb{Q}$-linear span
 of the operational classes associated to 
 all $[\Gamma_\delta,\gamma]\in \D\G_{g,n,d}$.
\end{definition}

By standard analysis \cite{GraberPandharipande}, the tautological classes
have a natural ring structure.
Our formula for $\mathsf{DR}_{g,A}^{\mathsf{op}}$ 
will be a sum of operational classes
determined by decorated prestable graphs of degree $d=\sum_{i=1}^n a_i$
(and hence will be tautological).

\subsubsection{Weightings mod \texorpdfstring{$r$}{r}} \label{Ssec:weightings}
\source{pixtons-formula-abel-jacobi-picard-stack}
Let $\Gamma_\delta\in \G_{g,n,d}$ be a prestable graph of degree $d$, and
let $r$ be a positive integer.

\begin{definition}
\source{pixtons-formula-abel-jacobi-picard-stack}
\notready A {\em weighting mod~$r$} of $\Gamma_\delta$ is a function on the set of half-edges,
$$ w:\H(\Gamma_\delta) \rightarrow \{0,1, \ldots, r-1\}\, ,$$
which satisfies the following three properties:

\begin{enumerate}
\item[(i)] $\forall i\in \L(\Gamma_\delta)$, corresponding to
 the marking $i\in \{1,\ldots, n\}$,
$$w(i)=a_i  \mod r\, ,$$
\item[(ii)] $\forall e \in \E(\Gamma_\delta)$, corresponding to two half-edges
$h,h' \in \H(\Gamma_\delta)$,
$$w(h)+w(h')=0 \mod r\, ,$$
\item[(iii)] $\forall v\in \V(\Gamma_\delta)$,
$$\sum_{v(h)= v} w(h)= \delta(v) \mod r\, ,$$ 
where the sum is taken over {\em all} $\mathsf{n}(v)$ half-edges incident 
to $v$.\end{enumerate}
\end{definition}



We denote by $\mathsf{W}_{\Gamma_\delta,r}$ the finite set of all possible weightings
mod $r$
of
$\Gamma_\delta$. The set $\mathsf{W}_{\Gamma_\delta,r}$ has cardinality $r^{h^1(\Gamma_\delta)}$.
We view $r$ as a {\em regularization parameter}.

\subsubsection{Calculation of the twisted double ramification cycle} \label{Ssec:MainFormula}
\source{pixtons-formula-abel-jacobi-picard-stack}
We denote by
$\P_{g,A,d}^{c,r}\in \mathsf{CH}^c_{\mathsf{op}}(\Picabs_{g,n,d})$
the codimension{\footnote{Codimension here is usually
called degree. But since we already have line bundle degrees,
we use the term codimension for clarity.}}
 $c$ component of the tautological operational class 
\begin{multline*}
\hspace{-10pt}\sum_{
\substack{\Gamma_\delta\in \G_{g,n,d} \\
w\in \mathsf{W}_{\Gamma_\delta,r}}
}
\frac{r^{-h^1(\Gamma_\delta)}}{|\Aut(\Gamma_\delta)| }
\;
j_{\Gamma_\delta*}\Bigg[
\prod_{i=1}^n \exp\left(\frac12 a_i^2 \psi_i + a_i \xi_i \right)
\prod_{v \in \V(\Gamma_\delta)} \exp\left(-\frac12 \eta(v) \right)
\\ \hspace{+10pt}
\prod_{e=(h,h')\in \E(\Gamma_\delta)}
\frac{1-\exp\left(-\frac{w(h)w(h')}2(\psi_h+\psi_{h'})\right)}{\psi_h + \psi_{h'}} \Bigg]\, .
\end{multline*} 
Several remarks about the formula are required:
\begin{enumerate}
    \item[(i)]
    The sum is over all isomorphism classes of prestable graphs of degree $d$ in the
    set $\mathsf{G}_{g,n,d}$. 
    Only finitely many underlying prestable graphs $\Gamma\in \mathsf{G}_{g,n}$ can
    contribute in fixed codimension $c$. However, for each such prestable graph, the above formula
    has infinitely many summands corresponding to the infinitely many functions
    $$\delta: \V \rightarrow \mathbb{Z}$$
    which satisfy the degree condition.
    The operational Chow class $\P_{g,A,d}^{c,r}$ is nevertheless well-defined since
      only finitely many summands have nonvanishing operation on any given family of curves
       carrying a degree $d$ line bundle over a base $\mathcal{S}$ of finite type.
    \item[(ii)] 
    Once the prestable graph $\Gamma_\delta$ is chosen, we sum over all $r^{h^1(\Gamma_\delta)}$ different 
    weightings $w\in \mathsf{W}_{\Gamma_\delta,r}$.
    \item[(iii)]
Inside the pushforward in the above formula, the first product 
$$\prod_{i=1}^n \exp\left(\frac{1}{2}a_i^2 \psi_{h_i}+a_i \xi_{h_i}\right)\, $$
is over $h\in \L(\Gamma)$
via the correspondence of legs and markings.
\item[(iv)]The class $\eta(v)$ is the $\eta_{0,2}$ class 
of Definition \ref{Not:PsiXiEta} associated to the vertex.
\item[(v)]
The third product is over all $e\in \E(\Gamma_\delta)$.
The factor 
$$\frac{1-\exp\left(-\frac{w(h)w(h')}{2}(\psi_h+\psi_{h'})\right)}{\psi_h + \psi_{h'}}$$
is well-defined since 
\begin{enumerate}
\item[$\bullet$]
the denominator formally divides
the numerator,
\item[$\bullet$] the factor is symmetric in $h$ and $h'$.
\end{enumerate}
No edge orientation is necessary.
\end{enumerate}

The following fundamental polynomiality property of $\P_{g,A,d}^{c,r}$ 
 is
parallel to Pixton's polynomiality in \cite[Appendix]{Janda2016Double-ramifica} and
is a consequence of \cite[Proposition $3''$]{Janda2016Double-ramifica}.

\begin{proposition}
\source{pixtons-formula-abel-jacobi-picard-stack}
\notready \label{pply} For fixed $g$, $A$, $d$, and $c$ and a decorated graph $[\Gamma_\delta, \gamma]$ of
\source{pixtons-formula-abel-jacobi-picard-stack}
degree $d$, the coefficient of $j_{\Gamma_\delta*} [\gamma]$ in the
tautological class
$$\P_{g,A,d}^{c,r} \in \mathsf{CH}^c_{\mathsf{op}}(\Picabs_{g,n,d})$$
is polynomial in $r$ (for all sufficiently large $r$).
\end{proposition}

We denote by $\P_{g,A,d}^c$ the value at $r=0$ 
of the polynomial associated to $\P_{g,A,d}^{c,r}$ by Proposition~\ref{pply}. In other words, $\P_{g,A,d}^c$ is the {\em constant} term of the associated polynomial in $r$. 

The main result of the paper is a 
 formula for the universal twisted double ramification cycle in operational Chow.{\footnote{Our
handling of the prefactor $2^{-g}$ in \cite[Theorem 1]{Janda2016Double-ramifica} differs here.
The factors of $2$ are now placed in the definition of $\P_{g,A,d}^{c,r}$
as in \cite{Janda2018Double-ramifica}}


\begin{theorem}
\source{pixtons-formula-abel-jacobi-picard-stack}
\notready \label{Thm:main}
\source{pixtons-formula-abel-jacobi-picard-stack}
Let $g\geq 0$ and $d\in \mathbb{Z}$.
Let $A=(a_1,\ldots,a_n)$ be a vector of integers with
$\sum_{i=1}^n a_i=d\,$.
The universal twisted double ramification cycle is calculated by Pixton's formula:
$$\mathsf{DR}^{\mathsf{op}}_{g,A} =\P_{g,A,d}^g\
\in {\mathsf{CH}}^g_{\mathsf{op}}(\Picabs_{g,n,d})\, .$$
\end{theorem}



Theorem \ref{Thm:main} is the most fundamental formulation of the relationship
between Abel-Jacobi theory and Pixton's formula that we know. Certainly,  
Theorem \ref{Thm:main} 
implies the double ramification cycle and $X$-valued
double ramification cycle results of \cite{Janda2016Double-ramifica,Janda2018Double-ramifica} }. But since we will
use \cite{Janda2018Double-ramifica} in the proof of Theorem \ref{Thm:main}, we provide no new approach to
these older results. However, the additional depth of Theorem \ref{Thm:main} immediately
allows new applications.


\subsection{Vanishing}
From his original double
ramification cycle formula,
Pixton conjectured an associated
vanishing property in the
tautological ring of the moduli
space of curves which
was proven by Clader and Janda \cite{cj}.
The parallel vanishing statement
in the tautological ring of the
moduli space of stable maps to $X$
was proven in \cite{Bae}.
The most general vanishing
statement is the following
result.


\begin{theorem}
\source{pixtons-formula-abel-jacobi-picard-stack}
\notready \label{Thm:mainvan}
\source{pixtons-formula-abel-jacobi-picard-stack}
Let $g\geq 0$ and $d\in \mathbb{Z}$.
Let $A=(a_1,\ldots,a_n)$ be a vector of integers with
$\sum_{i=1}^n a_i=d$.
 Pixton's vanishing holds
 in operational Chow:
$$ \P_{g,A,d}^c\ = 0 \
\in {\mathsf{CH}}^c_{\mathsf{op}}(\Picabs_{g,n,d})\, \ \ \ 
{\text{for all}} \ \ c>g\, .
$$ 
\end{theorem}

 



Theorem \ref{Thm:mainvan} may be viewed
as providing relations among
tautological classes in the operational
Chow ring of the Picard stack -- a new
direction of study with many open questions.{\footnote{See \cite{RandalWilliams2012,Yin}
for tautological relations on the Picard variety over
the moduli space of smooth curves.}}
While Theorem \ref{Thm:mainvan} 
implies the
vanishings of
\cite{Bae, cj}, 
we will use these results in our
proof. 


\subsection{Twisted holomorphic and meromorphic differentials} \label{Sect:Twistdiffintro}
\source{pixtons-formula-abel-jacobi-picard-stack}
\subsubsection{Fundamental classes}
Let $A=(a_1,\ldots,a_n)$ be a vector of zero and pole multiplicities
satisfying
$$\sum_{i=1}^n a_i= 2g-2\, .$$
Let $\HH_g(A) \subset \mathcal{M}_{g,n}$ be the quasi-projective locus 
of pointed curves $(C,p_1,\ldots,p_n)$ satisfying the condition
$$\mathcal{O}_C\Big(\sum_{i=1}^n a_i p_i\Big) \simeq \omega_C\, .$$
In other words, $\HH_g(A)$ is the locus of meromorphic differentials{\footnote{If all
the parts of $A$ are non-negative, then $\HH_g(A)$ is the locus of
holomorphic differentials.}}
with zero and pole multiplicities prescribed by $A$.
In \cite{Farkas2016The-moduli-spac}, a compact moduli space of twisted canonical divisors
$$\widetilde{\HH}_g(A) \subset \overline{\mathcal{M}}_{g,n}$$
is constructed which contains $\HH_g(A)$ as an open set.

In the strictly meromorphic
case, where $A$ contains at least one strictly negative part,
$\widetilde{\HH}_g(A)$ is of pure codimension $g$ in
$\overline{\mathcal{M}}_{g,n}$ by \cite[Theorem 3]{Farkas2016The-moduli-spac}. 
A weighted fundamental cycle of $\widetilde{\HH}_g(A)$, 
\begin{equation}\label{wwww4}
\source{pixtons-formula-abel-jacobi-picard-stack}
\mathsf{H}_{g,A}\in \mathsf{CH}_{2g-3+n}(\oM_{g,n})\ ,
\end{equation}
is constructed in
\cite[Appendix A]{Farkas2016The-moduli-spac} with explicit nontrivial 
weights on the irreducible components.
In the strictly meromorphic case,
$\HH_{g}(A)\subset \oM_{g,n}$ 
is also of pure codimension $g$.
The closure
$$\overline{\HH}_{g}(A)\subset \oM_{g,n}$$
contributes to the fundamental class $\mathsf{H}_{g,A}$ with
multiplicity 1, but there are additional boundary contributions, see
\cite[Appendix A]{Farkas2016The-moduli-spac}.

The  
universal family over the moduli space of
stable curves together with the 
relative dualizing sheaf,
\begin{equation}\label{rrrr}
\source{pixtons-formula-abel-jacobi-picard-stack}
\pi: \mathcal{C}_{g,n} \rightarrow \overline{\cM}_{g,n}\, , \ \ \ \ \omega_\pi \rightarrow \mathcal{C}_{g,n}\, ,
\end{equation}
determine an object of
$\Picabs_{g,n,2g-2}$.
By \cite{Holmes2019Infinitesimal-s} and
the compatibility of Theorem \ref{cccc},
the action of $\mathsf{DR}^{\mathsf{op}}_{g,A}$ on the fundamental class
of $\overline{\mathcal{M}}_{g,n}$ equals the  weighted
fundamental class of $\widetilde{\HH}_g(A)$,
$$\mathsf{DR}^{\mathsf{op}}_{g,A}(\phi_\omega)\Big( [\overline{\mathcal{M}}_{g,n}]\Big)= \mathsf{H}_{g,A}
\in {\mathsf{CH}}_{2g-3+n}(\overline{\mathcal{M}}_{g,n})
\, .$$
We can now apply Theorem \ref{Thm:main} to prove the following result.
 
\begin{theorem}
\source{pixtons-formula-abel-jacobi-picard-stack}
\notready\label{FPA1}
\source{pixtons-formula-abel-jacobi-picard-stack}
In the strictly meromorphic case,
$$\mathsf{H}_{g,A} = \mathsf{P}_{g,A,2g-2}^g[\oM_{g,n}]$$
for the universal family
$$\pi: \mathcal{C}_{g,n} \rightarrow \overline{\cM}_{g,n}\, ,\ \ \ \
\omega_\pi \rightarrow \mathcal{C}_{g,n}\,. $$
\end{theorem}


Theorem \ref{FPA1} is exactly equivalent to Conjecture A of 
\cite[Appendix A]{Farkas2016The-moduli-spac}.
Since both the moduli space  $\widetilde{\HH}_g(A)$ and the weighted
fundamental cycle 
$\mathsf{H}_{g,A}$
have explicit geometric definitions, 
the result provides a geometric representative
of Pixton's cycle class in terms of twisted differentials. 
%A different 
%approach to Conjecture A should be possibile
%using new logarithmic compactifications of the moduli of
%meromorphic differentials and the virtual localization formula
%in the log context.
Theorem \ref{FPA1} is proven in Section \ref{Sect:ProofConjA} where the parallel conjectures \cite{Schmitt2016Dimension-theor}
for higher differentials  are also proven (by parallel  arguments).

\subsubsection{Closures}
Let $A=(a_1,\ldots, a_n)$ be a vector of integers satisfying
$$\sum_{i=1}^n a_i =2g-2\, .$$
A careful investigation of the closure
$$\HH_g(A) \subset \overline{\HH}_g(A) \subset \oM_{g,n}$$
is carried out in \cite{Bainbridge2018Compactificatio} in
both the holomorphic and meromorphic cases.
By a simple procedure presented in \cite[Appendix]{Farkas2016The-moduli-spac},
Theorem \ref{FPA1} determines the cycle classes of the 
closures
$$ [\overline{\HH}_g(A)] \in \mathsf{CH}_{*} (\oM_{g,n})\, .$$
for $A$ in both the holomorphic and meromorphic cases. 

A similar procedure determines the corresponding classes for $k$-differentials, see \cite[Section 3.4]{Schmitt2016Dimension-theor} for an explanation. In particular, 
our results imply that the cycle classes of the closures 
are tautological{\footnote{The precise statement is given in
Corollary \ref{kk333} of Section \ref{clozzz}.}} 
for all $k$ (as was previously known only for $k=1$ due to \cite{Sauvaget2017Cohomology-clas}).

In the case of holomorphic differentials, another approach to the class of
the closure $\overline{\HH}_g(A) \subset \oM_{g,n}$ is provided
by Conjecture A.1 of the Appendix of \cite{Pandharipande2019Tautological-re} via a limit of
Witten's $r$-spin class. A significant first step in the proof of \cite[Conjecture A.1]{Pandharipande2019Tautological-re} by 
Chen, Janda, Ruan, and Sauvaget can be found in \cite{Chen2018Towards-a-Theor}. Further progress
requires 
a virtual localization analysis for moduli spaces of 
stable log maps. An approach to 
Theorem \ref{FPA1} using log stable maps, virtual localization in the log context, and
the strategy of \cite{Janda2016Double-ramifica} also appears possible (once the
required moduli spaces and localization formulas are established).


%\jocomment{In fact, Conjecture A for the $k$-differentials (in the case where one element of the partition is negative or not divisible by $k$ are the first proof that the corresponding cycles $[\overline{\HH}^k_g(A)]$ are tautological. In Adrien's paper \cite{Sauvaget2017Cohomology-clas} he only shows this for $k=1$. We should probably cite  his paper anyway.}
%\{Added a line to this effect above. Feel free to delete/edit, or remove  these comments if happy. }

\subsection{Invariance properties} \label{Sect:introinvariance}
\source{pixtons-formula-abel-jacobi-picard-stack}

%\Dcomment{Temporary paragraph proposing names for invariances}
%\begin{enumerate}
%    \item inverses \jocomment{alternative: dualizing} \Rcomment{I like "dualizing" invariance}
%    \item adding an unweighted marking\Rcomment{How about  "unweighted marking" invariance}
%    \item moving weights to the bundle \Rcomment{perhaps  "weight translation" invariance}
%    \item pullback from the base \Rcomment{perhaps  "pullback twisting" invariance}
%    \item twisting by a vertical divisor \Rcomment{perhaps  "vertical twisting" invariance}
%    \item subdividing the curve \jocomment{alternative (going the other way): partial stabilization of the curve}\Dcomment{Maybe J's is %better - the subdivision language is more logarithmic. }
%    \Rcomment{"partial stabilization" invariance}
%\end{enumerate}
%\jocomment{Is the idea of the referee to give a name to each invariance once the first time it is introduced, but continuing to use the %shorthands I - VI in the text, or writing the full name for each mention? If it is the second (which would be rather long and clumsy in the %main text) we could compromise and choose a 1-2 letter abbreviation for each, e.g. Invariance D(uality), Invariance U(nweighted marking), %Invariance W(eights to the bundle) etc.}
%\Dcomment{I guess the ref just wants it to be easier to rea/remember, I thik it's up to us what we do exactly. Your D, U, ... suggestion %seems fun to me,. }
%\Rcomment{names should have at most two words if we are going to use them again}.


The 
universal twisted double ramification cycle has several
basic invariance properties 
which play 
an important role in our study. 

Recall that an object of $\Picabs_{g,n,d}$ over $\mathcal{S}$ is
a flat family  of prestable
    $n$-pointed genus $g$ curves together with a line bundle
    of relative degree $d$,
    \begin{equation}
    \label{ff449911}    
\source{pixtons-formula-abel-jacobi-picard-stack}
    \pi:\mathcal{C} \rightarrow \mathcal{S}\, , \ \ \ \ p_1,\ldots,p_n: \mathcal{S} \rightarrow \mathcal{C}\, , \ \ \ \
        \mathcal{L} \rightarrow \mathcal{C}\, .
        \end{equation}
    Let $\mathsf{DR}^{\mathsf{op}}_{g,A,\mathcal{L}}\in \mathsf{CH}^g_{\mathsf{op}}(\mathcal{S})$
    be the twisted double ramification cycle associated to the above
    family \eqref{ff449911} and
     the vector
    $$A=(a_1,\ldots, a_n)\, , \ \ \ \ d =\sum_{i=1}^n a_i\, .$$


 
        \noindent{\bf{Invariance \invnrI: {\rm Dualizing}.}}
\vspace{10pt}

A new object of $\Picabs_{g,n,-d}$  over $\mathcal{S}$ is
obtained from \eqref{ff449911} by dualizing $\mathcal{L}$:
    \begin{equation}
    \label{ff44992}    
\source{pixtons-formula-abel-jacobi-picard-stack}
    \pi:\mathcal{C} \rightarrow \mathcal{S}\, , \ \ \ \ p_1,\ldots,p_n: \mathcal{S} \rightarrow \mathcal{C}\, , \ \ \ \
        \mathcal{L}^* \rightarrow \mathcal{C}\, .
        \end{equation}
 Let $\mathsf{DR}^{\mathsf{op}}_{g,-A, \mathcal{L}^*}\in \mathsf{CH}^g_{\mathsf{op}}(\mathcal{S})$ be the twisted double ramification cycle associated to the new family \eqref{ff44992} and the vector $-A=(-a_1,\ldots,-a_n)$. We have the invariance
        $$\mathsf{DR}^{\mathsf{op}}_{g,-A, \mathcal{L}^*}
        =\epsilon^* \mathsf{DR}^{\mathsf{op}}_{g,A, \mathcal{L}}
        \, , $$
where $\epsilon: \Picabs_{g,n,-d} \rightarrow \Picabs_{g,n,d}$
is the natural map obtained via dualizing the line bundle.
    

 \vspace{10pt}
        \noindent{\bf{Invariance \invnrn: {\rm Unweighted markings}.}}
\vspace{10pt}

Assume we have an additional section $p_{n+1}: \mathcal{S} \rightarrow \mathcal{C}$ of $\pi$ which yields an object of $\Picabs_{g,n+1,d}$,
    \begin{equation}
    \label{ff44993}    
\source{pixtons-formula-abel-jacobi-picard-stack}
    \pi:\mathcal{C} \rightarrow \mathcal{S}\, , \ \ \ \ p_1,\ldots,p_n,p_{n+1}: \mathcal{S} \rightarrow \mathcal{C}\, , \ \ \ \
        \mathcal{L} \rightarrow \mathcal{C}\, .
        \end{equation}
 Let $A_0 \in \mathbb{Z}^{n+1}$ be the vector obtained by appending $0$ (as the last coefficient) to $A$.
 Let $\mathsf{DR}^{\mathsf{op}}_{g,A_0, \mathcal{L}}\in \mathsf{CH}^g_{\mathsf{op}}(\mathcal{S})$ be the twisted double ramification cycle associated to the new family \eqref{ff44993} and the vector $A_0$. We have the invariance
        $$\mathsf{DR}^{\mathsf{op}}_{g,A_0, \mathcal{L}}
        =F^*\mathsf{DR}^{\mathsf{op}}_{g,A, \mathcal{L}}
        \, , $$
    where $F:\Picabs_{g,n+1,d} \rightarrow \Picabs_{g,n,d}$
    is the map forgetting the last marking.
 
 \vspace{10pt}
        \noindent{\bf{Invariance \invnrA: {\rm Weight translation}.}}
\vspace{10pt}

Let $B=(b_1,\ldots,b_n) \in \bb Z^n$ satisfy $\sum_{i=1}^n b_i=e$, then the family 
    \begin{equation}
    \label{ff44994}    
\source{pixtons-formula-abel-jacobi-picard-stack}
    \pi:\mathcal{C} \rightarrow \mathcal{S}\, , \ \ \ \ p_1,\ldots,p_n: \mathcal{S} \rightarrow \mathcal{C}\, , \ \ \ \
        \mathcal{L}\big(\sum_{i=1}^n b_i p_i\big) \rightarrow \mathcal{C}\, .
        \end{equation}
 defines an object of $\Picabs_{g,n,d+e}$. 
 Let $\mathsf{DR}^{\mathsf{op}}_{g,A+B, \mathcal{L}(\sum_i b_i p_i)}\in \mathsf{CH}^g_{\mathsf{op}}(\mathcal{S})$ be the twisted double ramification cycle associated to the new family \eqref{ff44994} and the vector $A+B$. We have the invariance
        $$\mathsf{DR}^{\mathsf{op}}_{g,A+B, \mathcal{L}(\sum_i b_i p_i)}
        =\mathsf{DR}^{\mathsf{op}}_{g,A, \mathcal{L}}
        \, . $$
    \vspace{10pt}   

        \noindent{\bf{Invariance \invnrII: {\rm Twisting by pullback.}}}
\vspace{10pt}
    
    Let $\mathcal{B} \rightarrow \mathcal{S}$ be any line bundle on the
    base. By tensoring \eqref{ff449911} with $\pi^*\mathcal{B}$, we obtain a new object of $\Picabs_{g,n,d}$ over $\mathcal{S}$:
    \begin{equation}
    \label{ff4499}    
\source{pixtons-formula-abel-jacobi-picard-stack}
    \pi:\mathcal{C} \rightarrow \mathcal{S}\, , \ \ \ \ p_1,\ldots,p_n: \mathcal{S} \rightarrow \mathcal{C}\, , \ \ \ \
        \mathcal{L}\otimes \pi^*\mathcal{B} \rightarrow \mathcal{C}\, .
        \end{equation}
        Let $\mathsf{DR}^{\mathsf{op}}_{g,A, \mathcal{L}\otimes\pi^*\mathcal{B}}\in \mathsf{CH}^g_{\mathsf{op}}(\mathcal{S})$ be the twisted double ramification cycle associated to the new family \eqref{ff4499} and the vector $A$. We have the invariance
        $$\mathsf{DR}^{\mathsf{op}}_{g,A, \mathcal{L}\otimes\pi^*\mathcal{B}}
        =\mathsf{DR}^{\mathsf{op}}_{g,A, \mathcal{L}}
        \, . $$
    
 \vspace{10pt}
        \noindent{\bf{Invariance \invnrIII: {\rm Vertical twisting}.}}
\vspace{10pt}
    
    Consider a partition of the genus, marking, and degree data,
    \begin{equation}\label{ss445}
\source{pixtons-formula-abel-jacobi-picard-stack}
    g_1+g_2=g\, , \ \ \ N_1\sqcup N_2 = \{1,\ldots,n\}\,, \ \ \ d_1+d_2=d\, ,
    \end{equation}
    which is {\em not} symmetric.{\footnote{We require
    $(g_1,N_1,d_1)\neq (g_2,N_2,d_2)$ so that the two sides of a separating node
    with separating data \eqref{ss445} can be distinguished.}}
    Such a partition defines a divisor 
    $$ \Delta_1 \in \mathsf{CH}^1(\mathfrak{C}_{g,n,d})$$
    in the universal curve over $\Picabs_{g,n,d}$
    by twisting by the $(g_1,N_1,d_1)$-component of a 
    curve with a separating node with separating data \eqref{ss445}.
 
    
    By tensoring \eqref{ff449911} with $\mathcal{O}_{\mathcal{C}}(\Delta_1)$, we obtain a new object of $\Picabs_{g,n,d}$ over $\mathcal{S}$:
    \begin{equation}
    \label{ff44999}    
\source{pixtons-formula-abel-jacobi-picard-stack}
    \pi:\mathcal{C} \rightarrow \mathcal{S}\, , \ \ \ \ p_1,\ldots,p_n: \mathcal{S} \rightarrow \mathcal{C}\, , \ \ \ \
        \mathcal{L}(\Delta_1) \rightarrow \mathcal{C}\, .
        \end{equation}
        Let $\mathsf{DR}^{\mathsf{op}}_{g,A, \mathcal{L}(\Delta_1)}\in \mathsf{CH}^g_{\mathsf{op}}(\mathcal{S})$ be the twisted double ramification cycle associated to the new family \eqref{ff44999} and the vector $A$. We have the invariance
        \begin{equation}\label{k55k9}
\source{pixtons-formula-abel-jacobi-picard-stack}
        \mathsf{DR}^{\mathsf{op}}_{g,A, \mathcal{L}(\Delta_1)} =
        \mathsf{DR}^{\mathsf{op}}_{g,A, \mathcal{L}}
        \, . \end{equation}
        
    For \emph{symmetric} separating data \eqref{ss445}, equality \eqref{k55k9}
    holds with $\Delta_1\subset \mathfrak{C}_{g,n,d}$ defined as the full preimage of the locus $\Delta \subset \Picabs_{g,n,d}$ of curves with a separating node \eqref{ss445}. Then, equality \eqref{k55k9}
    follows from Invariance \invnrII{} with $\mathcal{B}=\mathcal{O}(\Delta)$.
    
\vspace{10pt}
        \noindent{\bf{Invariance \invnrIV: {\rm Partial stabilization}.}}
\vspace{10pt}
    
    Consider a second family 
    of prestable
    $n$-pointed genus $g$ curves over $\mathcal{S}$,
    \begin{equation*}
    \pi':\mathcal{C}' \rightarrow \mathcal{S}\, , \ \ \ \ p'_1,\ldots,p'_n: \mathcal{S}
    \rightarrow \mathcal{C}'\, ,\end{equation*}
    together with a birational $\mathcal{S}$-morphism 
    $$f: \mathcal{C}' \rightarrow \mathcal{C}\, , \ \ \ \ f\circ p'_i = p_i\, .$$
    A line bundle of relative degree $d$ is defined on $\mathcal{C}'$ by
    $$f^*\mathcal{L} \rightarrow \mathcal{C}'\, .$$
    We require the following property to hold: {\em if the section $p'_i$ meets the
    exceptional locus of $f$, then $a_i=0$}. 
        
        Let $\mathsf{DR}^{\mathsf{op}}_{g,A, f^*\mathcal{L}}\in \mathsf{CH}^g_{\mathsf{op}}(\mathcal{S})$ be the twisted double ramification cycle associated to the new family
    \begin{equation} \label{eqn:primefamily}    \pi':\mathcal{C}' \rightarrow \mathcal{S}\, , \ \ \ \ p'_1,\ldots,p'_n: \mathcal{S}
\source{pixtons-formula-abel-jacobi-picard-stack}
    \rightarrow \mathcal{C}'\, , \ \ \ \ f^*\mathcal{L} \rightarrow \mathcal{C}' \end{equation}
    and the vector $A$.
        We have the invariance
        $$\mathsf{DR}^{\mathsf{op}}_{g,A, f^*\mathcal{L}}=
        \mathsf{DR}^{\mathsf{op}}_{g,A, \mathcal{L}}\, . $$


\vspace{10pt}
Theorem \ref{Thm:main} provides two paths to viewing 
the  above  invariance properties. The invariances can be 
seen
either from
formal
properties of the
geometric construction of the universal twisted double ramification
cycle or from formal symmetries of
Pixton's formula. In fact, all invariances hold not only for the codimension $g$ part $\P^g_{g,A,d}$ which computes the double ramification cycle, but for the full mixed-degree class $\P_{g,A,d}^\bullet$.

For example, Invariance \invnrIV{} on the 
formula side says that for the maps 
\[\phi_{\ca L}, \phi_{f^* \ca L} : \ca S \to \Picabs_{g,n,d}\]
obtained from the families  \eqref{ff449911} and \eqref{eqn:primefamily}, we have  
\[\phi_{f^* \ca L}^* \P^g_{g,A,d}=
\phi_{\ca L}^* \P^g_{g,A,d}  \in \mathsf{CH}^g_{\mathsf{op}}(\mathcal{S})\,  \] 
 for
$\P^g_{g,A,d} \in \CHop^g(\Picabs_{g,n,d})$.
%\in \prod_{c=0}^\infty \CHop^c(\Picabs_{g,n,d})\, .$$

Proofs of all of the invariances
will be presented in Section \ref{Sect:proofInvariance}.
The above invariances (together with geometric definitions
when transversality to the Abel-Jacobi map holds)  do not
characterize{\footnote {Further geometric properties are
required, see Section 1.6 of \cite{HolSch} for a discussion.}} 
$\DRop$. 

\subsection{Universal formula in degree 0}\label{sec:intro_deg_0}
\source{pixtons-formula-abel-jacobi-picard-stack}
The most efficient statement of the double ramification cycle formula
on the Picard stack of curves occurs in the degree $d=0$ case with {\em no} markings. 
In order to avoid{\footnote{For $g=1$, a parallel formula holds for $n=1$ and $A=(0)$.}} the unpointed genus 1 case, let $g\neq 1$.  


The specialization of Theorem \ref{Thm:main} to $d=0$ calculates $\DRop_{g,\emptyset}$ as
%The result $\P_g^g$ on $\Picabs_{g,0,0}=\Picabs_g$ is 
the value at $r=0$ of the degree $g$ part of
\begin{align*} %\label{eqn:P_gformula2}
\source{pixtons-formula-abel-jacobi-picard-stack}
\exp\left(-\frac12 \eta \right) \sum_{
\substack{\Gamma_\delta\in \G_{g,0,0} \\
w\in \mathsf{W}_{\Gamma_\delta,r}}
}
\frac{r^{-h^1(\Gamma_\delta)}}{|\Aut(\Gamma_\delta)| }
\;
j_{\Gamma_\delta*}\Bigg[&
% \prod_{v \in \V(\Gamma_\delta)} 
%\\ &
\prod_{e=(h,h')\in \E(\Gamma_\delta)}
\frac{1-\exp\left(-\frac{w(h)w(h')}2(\psi_h+\psi_{h'})\right)}{\psi_h + \psi_{h'}} \Bigg]\,  %\nonumber
\end{align*}
as an operational Chow class on 
$$\Picabs_g=\Picabs_{g,0,0}\, .$$
%As proven in  Lemma \ref{Lem:pullbackP_g} of 
The full
statement of Theorem \ref{Thm:main} can 
be recovered from the above $d=0$ specialization via pullback under the map
%that the full formula $\P_{g,A,d}^{g}$ can be reconstructed from $\P_g^g$ as a pullback under the maps
\[\Picabs_{g,n,d} \to \Picabs_{g}\, ,\ \ \  (C,p_1, \ldots, p_n, \ca L) \mapsto \left(C, \ca L\left(-\sum_{i=1}^n a_i p_i\right)\right).\]
Indeed, the above map is the composition of the morphism 
\[\tau_{-A} : \Picabs_{g,n,d} \to \Picabs_{g,n,0}, (C,p_1, \ldots, p_n, \ca L) \mapsto \left(C, p_1, \ldots, p_n, \ca L\left(-\sum_{i=1}^n a_i p_i\right)\right)\]
with the morphism $$F : \Picabs_{g,n,0} \to \Picabs_{g}$$ forgetting the markings $p_1, \ldots, p_n$. 

\vspace{5pt}
\noindent $\bullet$
For the $\DRop_{g,A}$ side of Theorem \ref{Thm:main}, 
Invariance \invnrn{} 
implies $$F^* \DRop_{g,\emptyset} = \DRop_{g,\bd 0}$$ for the zero vector $\bd 0 \in \mathbb{Z}^n$. Furthermore, Invariance \invnrA{} implies $$\tau_{-A}^*\DRop_{g,\bd 0} = \DRop_{g,A}\, .$$ 

\vspace{5pt}
\noindent $\bullet$ For the $\mathsf{P}^g_{g,A,d}$
side of Theorem \ref{Thm:main}, 
the corresponding invariance properties of Pixton's formula
(discussed in Section \ref{Sect:proofInvariance}) 
yield the parallel transformation 
$$\tau_{-A}^*F^*\mathsf{P}^g_{g,\emptyset,0} = \mathsf{P}^g_{g,A,d}\, .$$ 

\vspace{5pt}
\noindent Therefore, the equality in Theorem \ref{Thm:main} for general $A$ and $d$ follows from the specialization to $A=\emptyset$ and $d=0$.
For certain steps in our proof of Theorem \ref{Thm:main}, the 
$A=\emptyset$ and $d=0$
geometry is
advantageous and is used.

%Combining this with Theorems \ref{cccc} and \ref{Thm:main} we see that all classical double ramification cycles $\DR_{g,A,\omega^k}$ can be recovered from the single class $\P_g^g$ by pullbacks under maps
%\[\overline{\mathcal{M}}_{g,n} \to \Picabs_{g}, (C,p_1, \ldots, p_n) \mapsto \left(C, \omega_C^k\left(-\sum_{i=1}^n a_i p_i\right)\right).\]

By restricting to suitable open subsets of $\Picabs_{g}$, we can simplify the $d=0$ formula even further. Let $$\Picabs_{g}^\textup{ct} \subset \Picabs_{g}$$
be the locus where the curve $C$ is of compact type. 
We obtain 
\begin{equation}\label{ccttt}
\source{pixtons-formula-abel-jacobi-picard-stack}
 %\P_g^g |_{\Picabs_{g}^\textup{ct}} 
 \DRop_{g,\emptyset}|_{\Picabs_{g}^\textup{ct}} 
 = \frac{\theta^g}{g!} \,, \ \ \  
 \text{for}\ \ \theta = -\frac12 \left(\eta +  \sum_{\Delta} d_\Delta^2 [\Delta] \right),
\end{equation}
where the sum is over the boundary divisors $\Delta \subset \Picabs_g$ on which generically the curve splits into two components carrying line bundles of degrees $d_\Delta$ and $-d_\Delta$. 
The class $\theta$ here may be viewed as a universal theta divisor on $\Picabs_{g}^\textup{ct}$. 

Formula \eqref{ccttt} was first written on the moduli space of stable curves of compact type in 
\cite{Grushevsky2012The-double-rami,Hain2013Normal-function}. The operational Chow class  $\DRop_{g,\emptyset}$ on $\Picabs_g$,
however, is {\em not} the  power of a divisor. 

%\jocomment{Was this what you had in mind, Rahul? I suspect it would not be terribly hard to make the statements about $\theta$ above precise and prove them, but it needs some nontrivial argument, I think. We should probably also include some references to previous computations on compact type and the formula $[e]=\theta^g/g!$ above.}






\subsection{Acknowledgements}
We thank D. Chen, A. Chiodo, F. Janda, G. Farkas, T. Graber, S. Grushevsky, A. Kresch,
S. Molcho, M. M\"oller, A. Pixton, D. Ranganathan, A. Sauvaget, H.-H. Tseng,
J. Wise, and D. Zvonkine
for  
many discussions about Abel-Jacobi theory, double ramification cycles,
and meromorphic differentials. The AIM workshop on {\em Double ramification cycles and integrable systems} 
played a key role at the start of the paper. We are grateful to the organisers A. Buryak, R. Cavalieri, E. Clader, and P. Rossi. We thank the referees for several 
improvements in the presentation.



Y.~B. was supported by ERC-2017-AdG-786580-MACI and the
Korea Foundation for Advanced Studies.
D.~H. was partially supported by NWO grant 613.009.103. 
R.~P. was partially supported by 
SNF-200020-182181, SwissMAP, and
the Einstein Stiftung. 
J.~S. was supported by the SNF Early Postdoc.Mobility grant 184245 and thanks
the Max Planck Institute for Mathematics in Bonn for its hospitality. 
R.~S. was supported by  NWO grant 613.009.113. 


The project has received funding from the European Research
Council (ERC) under the European Union Horizon 2020 research and
innovation program (grant agreement No. 786580).
